\section{Problem Formulation}\label{sec:formulation}\index{Strategic feature revelation|textbf}\index{Feature withholding|textbf}

We model regression under strategic feature disclosure as a screening game between a sender (the feature provider) and a receiver (the decision maker). Let $\X \subseteq \R^d$ be the feature space and let the outcome space be $\Y := \R$. Let $(U,X,Y)$ be a tuple of random variables taking values in $\Y \times \X \times \Y$ with joint distribution $P_{(U,X,Y)}$. A realization $(u,x,y)$ represents a sender type with private features $x$, a preferred outcome $u$, and a ground-truth outcome $y$. The receiver predicts the outcome using a regression model $r: \X \cup \{\emptyset\} \to \Y$, and both players evaluate performance using the squared loss $\loss(\alpha,\beta) := (\alpha-\beta)^2$.

Nature draws a sender type $(u,x,y) \sim P_{(U,X,Y)}$. After observing its type, the sender chooses an action $s \in \{\emptyset, x\}$, where $\emptyset$ and $x$ represent withholding or revealing the feature vector, respectively. The chosen action passes through an erasure channel with parameter $q \in [0,1]$, arriving at the receiver as a random message $M \in \X \cup \{\emptyset\}$ with conditional distribution $P(m \mid s)$ given by
\begin{equation*}
    P(m \mid s) = 
    \begin{cases} 
    1 & \text{if } s = \emptyset \text{ and } m = \emptyset, \\
    1-q & \text{if } s = x \text{ and } m = x, \\
    q & \text{if } s = x \text{ and } m = \emptyset.
    \end{cases}
\end{equation*}
The receiver applies a regression rule $r$ to produce prediction $r(M)$. The sender incurs loss $\loss(r(M),u) = (u - r(M))^2$ and the receiver incurs loss $\loss(r(M),y) = (y - r(M))^2$.

The sender and receiver both seek to minimize their expected losses. The sender's strategy space is the collection of all functions mapping sender types to feasible actions, given by
\begin{equation*}
    \Sspace := \{s: \Y \times \X \to \X \cup \{\emptyset\} \mid s(u,x) \in \{\emptyset, x\} \text{ for all } (u,x) \in \Y \times \X\},
\end{equation*}
and the receiver's strategy space is\footnote{We restrict the receiver's strategy space to models without an intercept and where the default prediction is zero.}
\begin{equation*}
    \Rspace := \{r_b: \X \cup \{\emptyset\} \to \Y \mid r_b(\emptyset) = 0,\; r_b(x) = b^\top x \text{ for some } b \in \R^d\}.
\end{equation*}
In the screening game, the receiver first commits to a regression model $r_b \in \Rspace$. The sender, after observing its type $(u,x,y)$ and the receiver strategy $r_b$, best responds with
\begin{equation*}
    s_b(u,x) \in \arg\min_{a \in \{\emptyset, x\}} \E_{M \sim P(\cdot \mid a)}[\loss(r_b(M), u)] \quad \text{for all } (u,x) \in \Y \times \X.
\end{equation*}
The receiver anticipates this behavior and optimally chooses
\begin{equation*}
    b^\star \in \arg\min_{b \in \R^d} \E_{(U,X,Y)}\!\left[\E_{M \sim P(\cdot \mid s_b(U,X))}[\loss(r_b(M), Y)]\right].
\end{equation*}
The receiver-sender strategy pair $(r_{b^\star}, s_{b^\star})$ constitutes a Stackelberg equilibrium of the screening game.

We denote the environment by $\env := (U,X,Y,q,\Sspace,\Rspace)$. For an environment $\env$ and a receiver predictor $r_b \in \Rspace$, the vanilla risk is defined as
\begin{equation*}
    \RV(b) := q \E_{Y}[\loss(r_b(\emptyset), Y)] + (1-q) \E_{(X,Y)}[\loss(r_b(X), Y)].
\end{equation*}
This is the receiver's risk under the mandatory baseline (where the sender always reveals features, and omissions occur solely via channel erasure). Let $\RV^\star := \min_{b \in \R^d} \RV(b)$. For receiver predictor $r_b$ and sender best response $s_b$, the strategic risk is
\begin{equation*}
    \RS(b) := \E_{(U,X,Y)}\!\left[\E_{M \sim P(\cdot \mid s_b(U,X))}[\loss(r_b(M), Y)]\right],
\end{equation*}
which is the receiver's risk under voluntary feature disclosure. Let $\RS^\star := \min_{b \in \R^d} \RS(b)$. We define the strategic advantage of environment $\env$ as
\begin{equation*}
    \Delta(\env) := \RV^\star - \RS^\star.
\end{equation*}
A positive strategic advantage, $\Delta(\env) > 0$, indicates that voluntary feature revelation strictly improves receiver performance relative to the mandatory baseline.

Based on the receiver's knowledge of the environment $\env$, we analyze strategic regression under three distinct information structures. The first information structure is when the receiver has complete knowledge of the joint distribution $P_{(U,X,Y)}$, referred to as the \textit{Decision Problem}. The objective is to establish a necessary and a sufficient condition on the environment $\env$ and the receiver strategy space $\Rspace$ under which $\Delta(\env) > 0$.

The second information structure is when the receiver does not know $P_{(U,X,Y)}$, but observes a finite training dataset $D_n = \{(u_k, x_k, y_k)\}_{k=1}^n$ of i.i.d. samples, referred to as the \textit{Offline Learning Problem}. The objective is to learn $\Delta(\env)$ from the training data using an algorithm and derive uniform generalization bounds for it.

The third information structure is when the receiver does not know the environment and sequentially interacts with the unknown environment over $T$ rounds, observing only the realized bandit loss feedback after each action, referred to as the \textit{Online Learning Problem}. The objective is to learn the optimal disclosure regime and regression model from the interaction using an algorithm and derive regret bounds for it.
% Applied Fix: Replaced exact argmin with infima/approximate minimizers for strategic risk.
